\documentclass[11pt,a4paper,oneside]{book}
\usepackage{90-faelles/latex/studiearkiv}
\fancyhead[L]{\small\sffamily Fysik · statistik}
\hypersetup{pdftitle={Statistik til fysik: foreløbig kompendiestruktur},pdfsubject={Statistisk metode}}
\title{Statistik til fysik\\[4mm]\large Foreløbig kompendiestruktur}
\author{Fysik-studiearkiv}
\date{\today}
\begin{document}
\frontmatter
\maketitle
\chapter*{Grundlag og afgrænsning}
\addcontentsline{toc}{chapter}{Grundlag og afgrænsning}
Kompendiet udvikles selvstændigt på dansk med udgangspunkt i Roger Barlows
\emph{Statistics: A Guide to the Use of Statistical Methods in the Physical Sciences}.
Bogens udgave, indholdsfortegnelse og præcise notation er endnu ikke verificeret.
De følgende kapitler er foreløbige redaktionelle arbejdsområder, ikke en gengivelse
 af bogens kapitelrækkefølge. Brugerens præcise statistikpensum fastlægges senere.

Statistisk fysik hører fortsat til termodynamikfaget. Dette dokument behandler
statistisk metode. Se kildeposten STAT-BARLOW i repositoryets pensumregister.
\tableofcontents
\mainmatter
\chapter{Sandsynlighed og fordelinger}
\planlagt{Foreløbigt arbejdsområde; rækkefølge og notation afstemmes med Barlows verificerede indholdsfortegnelse. Pensumafgrænsning afventer.}
\section{Sandsynlighed, betingning og normalisering}
\leavevmode\par
% Tilføj selvstændig tekst, definitioner, antagelser og relevante figurer.

\section{Fordelinger, momenter og betingning}
\leavevmode\par
% Tilføj selvstændig tekst, definitioner, antagelser og relevante figurer.

\chapter{Estimatorer og estimation}
\planlagt{Foreløbigt arbejdsområde; rækkefølge og notation afstemmes med Barlows verificerede indholdsfortegnelse. Pensumafgrænsning afventer.}
\section{Estimatorers egenskaber og gyldighed}
\leavevmode\par
% Tilføj selvstændig tekst, definitioner, antagelser og relevante figurer.

\section{Estimation og konfidensintervaller}
\leavevmode\par
% Tilføj selvstændig tekst, definitioner, antagelser og relevante figurer.

\chapter{Likelihood}
\planlagt{Foreløbigt arbejdsområde; rækkefølge og notation afstemmes med Barlows verificerede indholdsfortegnelse. Pensumafgrænsning afventer.}
\section{Model, data og likelihoodfunktion}
\leavevmode\par
% Tilføj selvstændig tekst, definitioner, antagelser og relevante figurer.

\section{Parameterbestemmelse og usikkerhed}
\leavevmode\par
% Tilføj selvstændig tekst, definitioner, antagelser og relevante figurer.

\chapter{Hypotesetest}
\planlagt{Foreløbigt arbejdsområde; rækkefølge og notation afstemmes med Barlows verificerede indholdsfortegnelse. Pensumafgrænsning afventer.}
\section{Hypoteser og teststørrelser}
\leavevmode\par
% Tilføj selvstændig tekst, definitioner, antagelser og relevante figurer.

\section{Fejltyper og fortolkning}
\leavevmode\par
% Tilføj selvstændig tekst, definitioner, antagelser og relevante figurer.

\chapter{Regression og modelkontrol}
\planlagt{Foreløbigt arbejdsområde; rækkefølge og notation afstemmes med Barlows verificerede indholdsfortegnelse. Pensumafgrænsning afventer.}
\section{Regression og modelkontrol}
\leavevmode\par
% Tilføj selvstændig tekst, definitioner, antagelser og relevante figurer.

\section{Residualer og kontrol af modelantagelser}
\leavevmode\par
% Tilføj selvstændig tekst, definitioner, antagelser og relevante figurer.

\chapter{Usikkerhedsanalyse}
\planlagt{Foreløbigt arbejdsområde; rækkefølge og notation afstemmes med Barlows verificerede indholdsfortegnelse. Pensumafgrænsning afventer.}
\section{Usikkerhedspropagation og korrelation}
\leavevmode\par
% Tilføj selvstændig tekst, definitioner, antagelser og relevante figurer.

\section{Korrelationer og rapportering}
\leavevmode\par
% Tilføj selvstændig tekst, definitioner, antagelser og relevante figurer.

\end{document}
